\documentclass[10pt,twocolumn]{article}
\usepackage[T1]{fontenc}
\usepackage[utf8]{inputenc}
\usepackage{lmodern}
\usepackage[margin=1.8cm,top=2.4cm,bottom=2cm,columnsep=0.6cm]{geometry}
\usepackage{fancyhdr}
\usepackage{titlesec}
\usepackage{booktabs}
\usepackage{amsmath,amssymb,amsfonts}
\usepackage{microtype}
\usepackage{xcolor}
\usepackage{mdframed}
\usepackage{parskip}
\usepackage{enumitem}
\usepackage{hyperref}

\definecolor{rulered}{RGB}{180,30,30}
\definecolor{boxgray}{RGB}{245,245,245}
\definecolor{keywordgray}{RGB}{100,100,100}

\pagestyle{fancy}
\fancyhf{}
\fancyhead[L]{\textsc{\small Claude Mag}}
\fancyhead[C]{\small\textit{Machine Learning Research Digest}}
\fancyhead[R]{\small 2026-10-12}
\fancyfoot[C]{\small\thepage}
\renewcommand{\headrulewidth}{0.8pt}

\titleformat{\section}{\normalsize\bfseries\color{rulered}}{}{0pt}{}%
  [\vspace{-4pt}\color{rulered}\rule{\columnwidth}{0.4pt}]
\titlespacing{\section}{0pt}{8pt}{4pt}

\hypersetup{colorlinks=true,urlcolor=rulered,linkcolor=black}

\begin{document}
\setlength{\parindent}{0pt}
\setlength{\parskip}{4pt}

{\small\color{keywordgray}\textbf{Keywords:}
  diffusion models \textbullet{}
  pixel-space generation \textbullet{}
  text-to-image \textbullet{}
  VAE-free}\par\vspace{4pt}

{\LARGE\bfseries\raggedright
  Iris-3B: Going Beyond the Latent with Pixel-Space Diffusion
  Training, Conversion and Fine-Tuning}\par\vspace{4pt}

{\large\itshape\raggedright
  A 3B-Parameter VAE-Free Diffusion Model Matches Latent Baselines on
  Text-to-Image Quality but Finds No Advantage on Fine-Grained Pixel Tasks}\par\vspace{6pt}

{\small\textbf{By} Hanqiu Li Cai \& Chema Garabito (Speridlabs)}\par\vspace{2pt}

{\small\color{keywordgray}arXiv:2610.09450 \textbullet{} October 2026}
\par\vspace{2pt}\noindent{\color{rulered}\rule{\columnwidth}{0.6pt}}\vspace{6pt}

Iris-3B is a 3-billion-parameter text-to-image diffusion transformer trained
entirely in pixel space---no VAE encoder, no latent compression---via a
256$\to$512$\to$1024 resolution curriculum, achieving GenEval 0.798 and
DPG-Bench 86.52, matching latent FLUX.2 Klein. Fine-tuning on depth
estimation and super-resolution finds no improvement over the latent baseline,
suggesting the VAE bottleneck is less limiting than assumed.

\begin{mdframed}[backgroundcolor=boxgray,linecolor=rulered,linewidth=1pt,
    innerleftmargin=6pt,innerrightmargin=6pt,
    innertopmargin=6pt,innerbottommargin=6pt]
\textbf{\small AT A GLANCE}\par\vspace{4pt}
\begin{description}[leftmargin=0pt,labelsep=4pt,itemsep=2pt,parsep=0pt,
    font=\small\bfseries]
  \item[Problem:] VAE bottleneck in latent diffusion may limit pixel-level
    task quality; direct pixel-space diffusion is a natural alternative.
  \item[Why it matters:] If pixel-space diffusion is competitive,
    the field can revisit the latent paradigm assumption.
  \item[Method:] PiT pixel-space diffusion transformer, 3B parameters,
    trained with a 256$\to$512$\to$1024 resolution curriculum.
  \item[Result:] GenEval 0.798, DPG-Bench 86.52; no gain on depth
    estimation or 4$\times$ super-resolution over latent baselines.
  \item[Evidence:] Compared against latent FLUX.2 Klein fine-tunes on
    downstream tasks; converted pixel model underperforms both.
\end{description}
\end{mdframed}

\section{The Big Picture}

The dominant paradigm in text-to-image generation encodes images into a
latent space with a VAE before denoising with a diffusion model. This
reduces compute: a 1024$\times$1024 image is processed as a
$\sim$128$\times$128 latent, cutting memory and FLOPs by up to 64$\times$.
But the VAE reconstruction introduces artefacts---blurring, ringing, and
colour shift---that may degrade performance on pixel-level tasks such as
depth estimation and super-resolution.

Iris-3B asks: does removing the VAE bottleneck and training directly in
pixel space yield measurable gains on these tasks? The answer, after a
thorough empirical investigation, is: not yet.

\section{Why Existing Approaches Fall Short}

\textbf{Latent diffusion models} (LDMs) incur VAE reconstruction error at
inference time. For generation quality benchmarks, this error is tolerated;
for fine-grained pixel tasks, it may be a genuine bottleneck.

\textbf{Pixel-space diffusion at high resolution} is expensive: the denoising
network must process full-resolution feature maps, scaling quadratically
in spatial dimensions.

\textbf{Cascaded super-resolution} (used in early pixel-space models like
Imagen) adds architectural complexity and requires training separate upsampling
stages.

\textbf{Latent-to-pixel conversion} (fine-tuning a pre-trained latent model
in pixel space) is unexplored as an alternative to training from scratch.

\section{Core Mechanism}

\textbf{Iris-3B (from scratch):} A PiT (Pixel-space Image Transformer)
diffusion model with 3B parameters. Training uses a progressive resolution
curriculum: stage~1 at 256$^2$, stage~2 at 512$^2$, stage~3 at 1024$^2$.
Text conditioning uses Qwen3-VL-4B-Instruct as the text encoder.

\textbf{Converted pixel model:} FLUX.2 Klein 4B is adapted to pixel space
by replacing the VAE encoder/decoder with an identity mapping and adjusting
channel dimensions, then fine-tuned on the pixel-space denoising objective.

Both are then fine-tuned on downstream tasks (monocular depth, 4$\times$
super-resolution).

\section{Technical Formulation}

The pixel-space denoising objective is standard denoising score matching:
\[
\mathcal{L}_\text{DSM} = \mathbb{E}_{t,\mathbf{x}_0,\boldsymbol{\epsilon}}\!
\Bigl[\|\boldsymbol{\epsilon} - \boldsymbol{\epsilon}_\theta(\mathbf{x}_t, t, \mathbf{c})\|_2^2\Bigr]
\]
where $\mathbf{x}_t = \sqrt{\bar{\alpha}_t}\,\mathbf{x}_0 +
\sqrt{1-\bar{\alpha}_t}\,\boldsymbol{\epsilon}$ is the noisy image at timestep
$t$ and $\mathbf{c}$ is the text embedding.

The progressive training schedule uses a staged sequence
$(256^2, T_1) \to (512^2, T_2) \to (1024^2, T_3)$
with decreasing step counts at higher resolutions to avoid excessive
compute at the final stage.

For task-specific fine-tuning (e.g., depth estimation with ground-truth $\mathbf{d}$):
\[
\mathcal{L}_\text{task} = \mathcal{L}_\text{DSM}(\mathbf{d})
+ \lambda\|\hat{\mathbf{d}} - \mathbf{d}\|_1
\]

\section{Learning or Inference Procedure}

\textbf{Pre-training:}
\begin{enumerate}[itemsep=2pt,parsep=0pt]
  \item Stage 1: train at 256$^2$ on image--caption pairs for $T_1$ steps.
  \item Stage 2: upsample, continue at 512$^2$ for $T_2$ steps.
  \item Stage 3: continue at 1024$^2$ for $T_3$ steps.
\end{enumerate}

\textbf{Fine-tuning:} Replace conditioning input with the task signal;
train with the task-specific loss for a small number of steps.

\textbf{Inference:} Standard DDPM/DDIM sampling; no VAE decoding.

\section{What the Guarantee Says}

No theoretical guarantees are claimed. The paper makes empirical claims on
GenEval and DPG-Bench (generation quality) and on monocular depth
and DIV2K super-resolution (fine-grained tasks).

\section{Experimental Findings}

\begin{itemize}[itemsep=2pt,parsep=0pt]
  \item \textbf{Generation quality:} Iris-3B achieves \textbf{GenEval 0.798}
    and \textbf{DPG-Bench 86.52} at 1024$^2$, matching Qwen-Image and
    competitive with latent FLUX.2 Klein.
  \item \textbf{Monocular depth:} Iris-3B performs at parity with a latent
    FLUX.2 Klein fine-tune; the converted pixel model performs worse.
  \item \textbf{4$\times$ super-resolution (DIV2K):} Neither Iris-3B nor
    the converted model outperforms the latent FLUX.2 Klein fine-tune.
  \item \textbf{Inference cost:} Pixel-space at 1024$^2$ is substantially
    more expensive per step than latent inference.
\end{itemize}

\section{Ablations and Interpretation}

\textbf{Curriculum stages:} Skipping the 256$^2$ stage degrades final
generation quality; low-frequency pre-training is important for convergence.

\textbf{Conversion quality:} The converted pixel model underperforms both
the from-scratch model and the original latent model, suggesting that the
latent representation carries useful inductive bias.

\textbf{Text encoder:} Replacing Qwen3-VL-4B with a smaller encoder
degrades text alignment (DPG) more than object-count metrics (GenEval).

\textbf{Interpretation:} The null result suggests that either the VAE
bottleneck is less harmful than assumed, that the latent representation
is beneficial beyond compression, or that pixel-space models need more
task-specific data to close the gap.

\section*{Reference}

Hanqiu Li Cai, Chema Garabito (Speridlabs).
\textbf{Iris-3B: Going Beyond the Latent with Pixel-Space Diffusion
Training, Conversion and Fine-Tuning}.
arXiv:2610.09450, October 2026.\\
\url{https://arxiv.org/abs/2610.09450}

\end{document}
